% Options for packages loaded elsewhere
\PassOptionsToPackage{unicode}{hyperref}
\PassOptionsToPackage{hyphens}{url}
\documentclass[
  spanish,
  11pt,
  titlepage]{article}
\usepackage{xcolor}
\usepackage[margin=2.5cm]{geometry}
\usepackage{amsmath,amssymb}
\setcounter{secnumdepth}{5}
\usepackage{iftex}
\ifPDFTeX
  \usepackage[T1]{fontenc}
  \usepackage[utf8]{inputenc}
  \usepackage{textcomp} % provide euro and other symbols
\else % if luatex or xetex
  \usepackage{unicode-math} % this also loads fontspec
  \defaultfontfeatures{Scale=MatchLowercase}
  \defaultfontfeatures[\rmfamily]{Ligatures=TeX,Scale=1}
\fi
\usepackage{lmodern}
\ifPDFTeX\else
  % xetex/luatex font selection
  \setmainfont[]{Times New Roman}
\fi
% Use upquote if available, for straight quotes in verbatim environments
\IfFileExists{upquote.sty}{\usepackage{upquote}}{}
\IfFileExists{microtype.sty}{% use microtype if available
  \usepackage[]{microtype}
  \UseMicrotypeSet[protrusion]{basicmath} % disable protrusion for tt fonts
}{}
\usepackage{setspace}
\makeatletter
\@ifundefined{KOMAClassName}{% if non-KOMA class
  \IfFileExists{parskip.sty}{%
    \usepackage{parskip}
  }{% else
    \setlength{\parindent}{0pt}
    \setlength{\parskip}{6pt plus 2pt minus 1pt}}
}{% if KOMA class
  \KOMAoptions{parskip=half}}
\makeatother
\usepackage{longtable,booktabs,array}
\usepackage{calc} % for calculating minipage widths
% Correct order of tables after \paragraph or \subparagraph
\usepackage{etoolbox}
\makeatletter
\patchcmd\longtable{\par}{\if@noskipsec\mbox{}\fi\par}{}{}
\makeatother
% Allow footnotes in longtable head/foot
\IfFileExists{footnotehyper.sty}{\usepackage{footnotehyper}}{\usepackage{footnote}}
\makesavenoteenv{longtable}
\usepackage{graphicx}
\makeatletter
\newsavebox\pandoc@box
\newcommand*\pandocbounded[1]{% scales image to fit in text height/width
  \sbox\pandoc@box{#1}%
  \Gscale@div\@tempa{\textheight}{\dimexpr\ht\pandoc@box+\dp\pandoc@box\relax}%
  \Gscale@div\@tempb{\linewidth}{\wd\pandoc@box}%
  \ifdim\@tempb\p@<\@tempa\p@\let\@tempa\@tempb\fi% select the smaller of both
  \ifdim\@tempa\p@<\p@\scalebox{\@tempa}{\usebox\pandoc@box}%
  \else\usebox{\pandoc@box}%
  \fi%
}
% Set default figure placement to htbp
\def\fps@figure{htbp}
\makeatother
\ifLuaTeX
\usepackage[bidi=basic]{babel}
\else
\usepackage[bidi=default]{babel}
\fi
\ifPDFTeX
\else
\babelfont{rm}[]{Times New Roman}
\fi
% get rid of language-specific shorthands (see #6817):
\let\LanguageShortHands\languageshorthands
\def\languageshorthands#1{}
\setlength{\emergencystretch}{3em} % prevent overfull lines
\providecommand{\tightlist}{%
  \setlength{\itemsep}{0pt}\setlength{\parskip}{0pt}}
\usepackage[style=apa,]{biblatex}
\addbibresource{referencias.bib}
\usepackage{float}
\usepackage{titlesec}
\titleformat{\section}{\fontsize{16}{18}\bfseries}{\thesection}{1em}{}
\titleformat{\subsection}{\fontsize{14}{16}\bfseries}{\thesubsection}{1em}{}
\definecolor{gris}{RGB}{89, 88, 88}
\usepackage{titling}
\pretitle{\begin{center}\vspace{\fill}\fontsize{16}{18}\selectfont}
\posttitle{\par\end{center}}
\preauthor{\begin{center}\fontsize{11}{13}\selectfont}
\postauthor{\par\end{center}}
\predate{\begin{center}\fontsize{11}{13}\selectfont}
\postdate{\end{center}\vspace{\fill}}
\usepackage{subfig}
\usepackage{bookmark}
\IfFileExists{xurl.sty}{\usepackage{xurl}}{} % add URL line breaks if available
\urlstyle{same}
\hypersetup{
  pdftitle={Trabajo Práctico 2},
  pdflang={es},
  hidelinks,
  pdfcreator={LaTeX via pandoc}}

\title{\textbf{Trabajo Práctico 2}}
\usepackage{etoolbox}
\makeatletter
\providecommand{\subtitle}[1]{% add subtitle to \maketitle
  \apptocmd{\@title}{\par {\large #1 \par}}{}{}
}
\makeatother
\subtitle{Laboratorio de Datos - Instituto de Cálculo\\
Facultad de Ciencias Exactas y Naturales - Universidad de Buenos Aires}
\author{Grupo Nº 3:\\
Gesualdi Santiago y Palmer Candela Magalí}
\date{2026-06-29}

\begin{document}
\maketitle

{
\setcounter{tocdepth}{3}
\tableofcontents
}
\setstretch{1.5}
\newpage

\section{Introducción}\label{introducciuxf3n}

El mercado inmobiliario es un sector dinámico cuya comprensión resulta fundamental para diversos actores económicos y sociales. En este contexto, el presente informe analiza un conjunto de datos que reúne información sobre precios de viviendas en la Ciudad Autónoma de Buenos Aires (CABA) y el Gran Buenos Aires correspondiente al año 2020. Los datos analizados provienen de la plataforma Kaggle, en particular a ``Buenos Aires Real State on sale''\autocite{properati} y han sido adaptados con fines didácticos para este estudio.
La importancia de abordar este análisis radica en la necesidad de identificar qué factores, ya sean características edilicias (como superficie, cantidad de ambientes o baños) o localización geográfica, inciden de manera determinante en la valoración de una propiedad. Comprender estas relaciones es clave para interpretar el funcionamiento del mercado, detectar patrones de comportamiento y aportar evidencia que permita tomar decisiones informadas en contextos de compra, venta o planificación urbana.
Además, el mercado inmobiliario presenta una marcada heterogeneidad, propiedades con configuraciones similares pueden exhibir precios muy distintos según su ubicación, distribución interna o condiciones edilicias. Esta variabilidad hace que el análisis estadístico y el modelado predictiva resulten herramientas especialmente valiosas para estudiar esta dinámica. En este trabajo se emplean métodos de regresión lineal simple y múltiple para evaluar la capacidad de distintas variables de explicar el precio de las viviendas, así como técnicas de agrupamiento para identificar conjuntos de propiedades con características estructurales comparables dentro de la región de Capital Federal.
El problema central que busca resolver este trabajo es la identificación de patrones y relaciones entre las variables físicas de las viviendas y su valor de mercado, así como la agrupación de las propiedades en conjuntos con características homogéneas. En función de ello, el objetivo principal del informe es evaluar la capacidad predictiva de distintos modelos y caracterizar los grupos resultantes, aportando una visión integral del comportamiento del mercado inmobiliario en el área analizada.

\section{Metodología}\label{metodologuxeda}

\begin{table}[H][!h]
\centering
\caption{(\#tab:descripcion_numericas)Estadísticos descriptivos de las variables númericas}
\centering
\begin{tabular}[t]{l|r|r|r|r|r|r}
\hline
\textbf{Variable} & \textbf{Media} & \textbf{Mediana} & \textbf{Desvío estándar} & \textbf{Mínimo} & \textbf{Máximo} & \textbf{Datos Faltantes}\\
\hline
bathrooms & 1.54 & 1.0 & 0.80 & 1 & 14 & 1,232\\
\hline
bedrooms & 2.03 & 2.0 & 1.05 & 0 & 15 & 0\\
\hline
price & 212,518.37 & 167,949.5 & 150,116.63 & 15,859 & 999,999 & 0\\
\hline
rooms & 3.07 & 3.0 & 1.30 & 1 & 16 & 0\\
\hline
surface\_covered & 87.09 & 66.0 & 73.94 & 11 & 5,250 & 0\\
\hline
surface\_total & 132.27 & 76.0 & 197.78 & 11 & 5,650 & 0\\
\hline
\end{tabular}
\end{table}

\begin{table}[H][!h]
\centering
\caption{(\#tab:descripcion_categoricas)Distribución de frecuencias de las variables categóricas principales}
\centering
\begin{tabular}[t]{l|l|r|r|r}
\hline
\textbf{Variable} & \textbf{Categoria} & \textbf{Frecuencia absoluta} & \textbf{Frecuencia relativa} & \textbf{Datos faltantes}\\
\hline
Region & Bs.As. G.B.A. Zona Norte & 15,937 & 0.18 & 0\\
\hline
Region & Bs.As. G.B.A. Zona Oeste & 6,297 & 0.07 & 0\\
\hline
Region & Bs.As. G.B.A. Zona Sur & 8,070 & 0.09 & 0\\
\hline
Region & Capital Federal & 59,124 & 0.66 & 0\\
\hline
Tipo de propiedad & Casa & 11,486 & 0.13 & 0\\
\hline
Tipo de propiedad & Departamento & 67,503 & 0.75 & 0\\
\hline
Tipo de propiedad & PH & 10,439 & 0.12 & 0\\
\hline
\end{tabular}
\end{table}

\begin{figure}[H]

{\centering \subfloat[Conjunto completo, con colas largas.(\#fig:surface_dist-1)]{\includegraphics[width=0.48\linewidth]{Resolución-TP-2_files/figure-latex/surface_dist-1} }\subfloat[Tras recortar el percentil 99 superior.(\#fig:surface_dist-2)]{\includegraphics[width=0.48\linewidth]{Resolución-TP-2_files/figure-latex/surface_dist-2} }

}

\caption{Distribución de la superficie cubierta de las propiedades (escala logarítmica).}(\#fig:surface_dist)
\end{figure}

\section{Resultados}\label{resultados}

\begin{figure}[H]

{\centering \subfloat[Escala lineal.(\#fig:price_dist-1)]{\includegraphics[width=0.48\linewidth]{Resolución-TP-2_files/figure-latex/price_dist-1} }\subfloat[Escala logarítmica.(\#fig:price_dist-2)]{\includegraphics[width=0.48\linewidth]{Resolución-TP-2_files/figure-latex/price_dist-2} }

}

\caption{Distribución del precio de venta de las propiedades.}(\#fig:price_dist)
\end{figure}

\begin{figure}[H]

{\centering \includegraphics[width=0.8\linewidth]{Resolución-TP-2_files/figure-latex/densidadesXregion-1} 

}

\caption{Distribución del precio mediano por m² según región.}\label{fig:densidadesXregion}
\end{figure}

\begin{figure}[H]

{\centering \includegraphics[width=0.95\linewidth]{Resolución-TP-2_files/figure-latex/precioXbarrio-1} 

}

\caption{Precio mediano por m² por barrio o localidad del Gran Buenos Aires, segmentado por zona. Las barras representan el rango intercuartílico.}\label{fig:precioXbarrio}
\end{figure}

\begin{figure}[H]

{\centering \subfloat[\label{fig:box-precio-1}]{\includegraphics[width=0.49\linewidth]{Resolución-TP-2_files/figure-latex/box-precio-1} }\subfloat[\label{fig:box-precio-2}]{\includegraphics[width=0.49\linewidth]{Resolución-TP-2_files/figure-latex/box-precio-2} }

}

\caption{Precio de venta según (a) tipo de propiedad y (b) región.}\label{fig:box-precio}
\end{figure}

En la matriz de correlacion se observa como las variables que mas se correlacionan con el precio son bathrooms, surface\_covered, rooms y bedrooms, asi que evaluaremos cada modelo con esas vars como predictoras de precio

\}
Para cumplir con el objetivo de maximizar la capacidad predictiva del modelo (Consigna 2), no se seleccionó el predictor basándose únicamente en el ajuste del set de entrenamiento (R2), el cual puede ser sensible a puntos de alta palanca. En su lugar, se compararon los modelos candidatos mediante Validación Cruzada de 10 folds, seleccionando la variable que minimizó el Error Cuadrático Medio de validación, garantizando así una mejor generalización ante datos no observados.

\begin{table}[H][!h]
\centering
\caption{(\#tab:rls_composicion)Modelos de RLS ordenados por capacidad explicativa}
\centering
\begin{tabular}[t]{l|r|r|r}
\hline
\textbf{Predictora} & \textbf{r2} & \textbf{MSE} & \textbf{RMSE}\\
\hline
bathrooms & 0.40 & 12,753,190,360 & 112,930.0\\
\hline
surface\_covered & 0.39 & 12,926,022,364 & 113,692.7\\
\hline
rooms & 0.25 & 15,926,934,113 & 126,202.0\\
\hline
bedrooms & 0.24 & 16,037,913,358 & 126,640.9\\
\hline
\end{tabular}
\end{table}

Si bien la variable bathrooms es la que menos RMSE presenta, al ser discreta y no continua, no construye un buen modelo predictivo para price por si sola, esto incluso lo podemos visualizar en el siguiente scatter plot donde se observa que claramente no hay una relacion linear y el modelo no ajusta bien. Por lo que, elegiremos surface\_covered como variable predictora

\begin{center}\includegraphics{Resolución-TP-2_files/figure-latex/rls_bathrooms-1} \end{center}

\begin{figure}[H]

{\centering \subfloat[Recta ajustada sobre los datos.(\#fig:rls_surface_plot-1)]{\includegraphics[width=0.48\linewidth]{Resolución-TP-2_files/figure-latex/rls_surface_plot-1} }\subfloat[Residuos frente a la superficie cubierta.(\#fig:rls_surface_plot-2)]{\includegraphics[width=0.48\linewidth]{Resolución-TP-2_files/figure-latex/rls_surface_plot-2} }

}

\caption{Regresión lineal simple del precio sobre la superficie cubierta (muestra aleatoria de 15.000 propiedades).}(\#fig:rls_surface_plot)
\end{figure}

\begin{table}[H][!h]
\centering
\caption{(\#tab:rlm_seleccion_variables)Comparación de modelos de RLM según variables predictoras utilizadas respectivamente.}
\centering
\begin{tabular}[t]{l|l|r|r}
\hline
\textbf{Modelo} & \textbf{Predictores (acumulados)} & \textbf{MSE validación (×10⁹)} & \textbf{MSE entrenamiento (×10⁹)}\\
\hline
M 1 & Baños & 12.75 & 12.75\\
\hline
M 2 & + Barrio/Localidad & 9.52 & 9.50\\
\hline
M 3 & + Superficie cubierta & 6.66 & 6.65\\
\hline
M 4 & + Tipo de propiedad & 6.12 & 6.10\\
\hline
M 5 & + Superficie total & 6.11 & 6.10\\
\hline
M 6 & + Ambientes & 6.10 & 6.09\\
\hline
M 7 & + Dormitorios & 6.10 & 6.08\\
\hline
M 8 & + Región & 6.10 & 6.08\\
\hline
\end{tabular}
\end{table}

Para determinar la cantidad óptima de variables en el modelo de regresión lineal múltiple, analizamos el MSE de validación obtenido mediante validación cruzada para modelos con 1 a 8 predictoras.
La curva de MSE muestra una disminución pronunciada entre 1 y 3 variables, seguida de una mejora moderada al incorporar la cuarta. A partir de ese punto, el MSE se estabiliza y las diferencias entre modelos más complejos son marginales (del orden de 0.02 × 10⁹), lo cual resulta irrelevante dada la escala de la variable respuesta.
Por criterios de parsimonia y estabilidad, seleccionamos el modelo con 4 variables, correspondiente al punto de inflexión (``codo'') de la curva. Este modelo logra un equilibrio adecuado entre capacidad predictiva y complejidad, evitando incorporar predictoras que no aportan mejoras significativas en la generalización.

\begin{figure}[H]

{\centering \includegraphics[width=0.75\linewidth]{Resolución-TP-2_files/figure-latex/rlm_codo-1} 

}

\caption{Error cuadrático medio de validación cruzada según el número de variables incluidas en el modelo de RLM.}(\#fig:rlm_codo)
\end{figure}

\begin{figure}[H]

{\centering \includegraphics[width=0.95\linewidth]{Resolución-TP-2_files/figure-latex/rlm_pred-1} 

}

\caption{Precios predichos por el modelo final de RLM frente a los precios reales observados.}(\#fig:rlm_pred)
\end{figure}

Para evitar problemas de memoria y acelerar el cálculo, se utilizó una muestra aleatoria de 5,000 observaciones, suficiente para capturar la estructura general del dataset y obtener una estimación estable del número óptimo de clústeres.

\begin{verbatim}
##  surface_covered        rooms             bedrooms          bathrooms       
##  Min.   :-1.35116   Min.   :-1.60349   Min.   :-1.95744   Min.   :-0.64791  
##  1st Qu.:-0.68111   1st Qu.:-0.77422   1st Qu.:-0.95253   1st Qu.:-0.64791  
##  Median :-0.29205   Median : 0.05504   Median : 0.05238   Median :-0.64791  
##  Mean   : 0.01956   Mean   : 0.01557   Mean   : 0.01399   Mean   : 0.01744  
##  3rd Qu.: 0.39961   3rd Qu.: 0.88431   3rd Qu.: 1.05729   3rd Qu.: 0.68386  
##  Max.   : 5.93289   Max.   : 5.85991   Max.   : 6.08185   Max.   :16.66510
\end{verbatim}

\begin{table}[H][!h]
\centering
\caption{\label{tab:barrios-caros-caba}Barrios de CABA más valorizados por precio del m² cubierto (mín. 200 registros).}
\centering
\begin{tabular}[t]{lrrr}
\toprule
\textbf{Barrio} & \textbf{N} & \textbf{Mediana del precio [USD]} & \textbf{Mediana del precio por m² [USD/m²]}\\
\midrule
Puerto Madero & 828 & 545,000 & 5,779\\
Las Cañitas & 465 & 340,000 & 4,000\\
Nuñez & 1,291 & 210,000 & 3,553\\
Palermo & 8,670 & 220,000 & 3,544\\
Belgrano & 4,292 & 250,000 & 3,515\\
\addlinespace
Recoleta & 3,898 & 285,045 & 3,383\\
Villa Urquiza & 2,183 & 175,000 & 3,182\\
Barrio Norte & 2,273 & 230,000 & 3,158\\
\bottomrule
\end{tabular}
\end{table}

\begin{figure}[H]

{\centering \includegraphics[width=0.95\linewidth]{Resolución-TP-2_files/figure-latex/elbow_kmeans-1} 

}

\caption{Método del codo: suma de cuadrados intra-clúster (WSS) en función del número de clústeres k.}(\#fig:elbow_kmeans)
\end{figure}

El codo se observa en k = 3, donde la disminución de WSS deja de ser pronunciada. Por encima de ese valor, las mejoras marginales son pequeñas, indicando que 3 clústeres capturan adecuadamente la estructura del conjunto de datos.

\begin{figure}[H]

{\centering \includegraphics[height=8cm]{Resolución-TP-2_files/figure-latex/cluster-viz-1} 

}

\caption{Segmentos de viviendas de CABA proyectados sobre las dos primeras componentes principales (muestra de 5.000 propiedades).}\label{fig:cluster-viz}
\end{figure}

\begin{verbatim}
## <ggplot2::labels> List of 1
##  $ title: NULL
\end{verbatim}

\begin{table}[H]
\centering
\caption{\label{tab:perfil}Caracterización de los 3 segmentos de viviendas de CABA (k-means).}
\centering
\fontsize{9}{11}\selectfont
\begin{tabular}[t]{lrrrrrr}
\toprule
\textbf{Segmento} & \textbf{N} & \textbf{Sup. cub. (m²)} & \textbf{Ambientes} & \textbf{Dormit.} & \textbf{Baños} & \textbf{Precio mediano [USD]}\\
\midrule
1 · Compactos & 23,325 & 42.65 & 1.8 & 1.0 & 1.1 & 120,000\\
2 · Estándar & 25,701 & 75.27 & 3.3 & 2.3 & 1.4 & 199,000\\
3 · Premium & 9,835 & 154.09 & 4.7 & 3.4 & 2.6 & 389,999\\
\bottomrule
\end{tabular}
\end{table}

\section{Discusiones}\label{discusiones}

\section{Referencias}\label{referencias}

\printbibliography[heading=none]

\renewcommand{\printbibliography}[1][]{}

\section{Anexos}\label{anexos}

\begin{center}\includegraphics{Resolución-TP-2_files/figure-latex/tablita_na-1} \end{center}

\begin{figure}[H]

{\centering \includegraphics[height=7.5cm]{Resolución-TP-2_files/figure-latex/cluster_distancias-1} 

}

\caption{Matriz de distancias en el espacio estandarizado (superficie cubierta, ambientes, dormitorios y baños). La diagonal (recuadrada) mide la cohesión intra-clúster como distancia media de las propiedades a su centroide; fuera de la diagonal, la separación inter-clúster entre centroides. Valores bajos en la diagonal y altos fuera de ella indican grupos compactos y bien separados.}(\#fig:cluster_distancias)
\end{figure}

\printbibliography

\end{document}
